\documentclass[../Main.tex]{subfiles}

\begin{document}

\section{Einleitung}

\subsection{Motivation und Relevanz}

Diese Maturitätsarbeit entstand aus dem Wunsch, Spiele nicht nur zu spielen, sondern sie genauer zu verstehen und selbst ein solches Spielerlebnis zu erschaffen.

Das Hack-and-Slash-Genre stellt aus unserer Sicht besonders hohe Anforderungen an das Zusammenspiel einzelner Spielsysteme, da Spielerbewegung, Kampfsystem, Gegnerverhalten und Leveldesign in Echtzeit ineinandergreifen. Im Zentrum unseres Interesses stehen das Gegnerverhalten und die Gegnerplatzierung. Wir haben diese beiden Bereiche gewählt, weil sie unserer Einschätzung nach für das Genre besonders tragend sind. Da das Hack-and-Slash-Genre auf Kampfgeschehen aufbaut, glauben wir, dass die Gegner besonders wichtig sind.

Man kann viel Forschung zu Videospielen finden. Wir haben auch im Rahmen dieser Maturitätsarbeit auf viel Recherche, wie z.b. "MDA", "GDC-Talks" oder "Game User Research", zurück gegriffen. Die genannten Beispiele werden in kapitel 2 Methodik dieser Arbeit erklärt. Was sie aber von uns unterscheidet ist, dass sie oft Spiele allgemein behandeln und nicht wie wir spezifisch das Hack-and-Slash-Genre behandeln.

In unserer Wahrnehmung erschienen Genrevertreter wie \game{Ninja Gaiden Sigma Plus} oder \game{NieR: Automata} in einer Phase, in der das Genre vergleichsweise häufig neue Titel hervorbrachte, während heute nur noch wenige grosse Titel erscheinen. Dadurch gibt es einerseits genügend Referenzmaterial, um zentrale Genremerkmale herauszuarbeiten. Andererseits sind wir bei der Umsetzung nicht an aktuelle Genrekonventionen gebunden, was uns zusätzlichen Freiraum für eigene, moderne Interpretationen klassischer Mechaniken lässt.

Unser Prototyp richtet sich in erster Linie an Fans der oben genannten Spiele. Da das Hack-and-Slash-Genre ein Untergenre der Action-Spiele ist, zählen wir auch Fans von Action-Spielen allgemein zu unserer Zielgruppe.

\begin{itemize}[leftmargin=1.5cm]
    \item \textcolor{red}{\textbf{Eingrenzung:} Die Begründung ist jetzt da, aber noch eine Behauptung. \emph{Warum} sind Gegnerverhalten und Gegnerplatzierung tragend? Ein Satz mit eurem Grund (z.\,B. was sie im Kampf für den Spieler bestimmen).}
    \item \textcolor{red}{\textbf{Kurzer Forschungsstand} (Checkliste) fehlt noch: Was gibt es bereits? MDA (Hunicke), Game Feel (Swink), Leveldesign-Literatur (Totten, Rogers, GDC), Games User Research (Drachen) -- und welche Lücke bleibt: Keine dieser Quellen behandelt das Hack-and-Slash-Genre direkt (vgl.\ Abschnitt~\ref{sec:VorgehenLeveldesign}). 2--4 Sätze, Details in Kapitel~\ref{sec:Grundlagen}.}
    \item \textcolor{red}{\textbf{Referenzspiele:} In Kapitel~\ref{sec:Referenzspielanalyse} analysiert ihr \game{Ninja Gaiden Sigma Plus} und \game{NieR: Automata}, nicht \game{God of War}. \game{NieR: Automata} hier auch nennen, sonst wundert sich der Leser. Bei \game{God of War} angeben, welcher Teil (2005 oder 2018). Alle genannten Spiele in die Ludografie.}
    \item \textcolor{red}{\enquote{heute nur noch wenige grosse Titel}: Durch \enquote{in unserer Wahrnehmung} ist das als Einschätzung markiert, das reicht. Falls ihr es belegen wollt: Erscheinungsjahre einiger Titel nennen.}
    \item \textcolor{red}{Optional für einen spannenden Einstieg: mit einer konkreten Spielszene beginnen (z.\,B. ein Gegner in \game{Ninja Gaiden}, der dich flankiert, während du einen anderen bekämpfst).}
\end{itemize}

\subsection{Erkenntnisfrage} \label{sec:Erkenntnisfrage}

Im Zentrum unserer Arbeit steht folgende Erkenntnisfrage:

\vspace{\baselineskip}

\begin{quote}
    \textit{Wie müssen Gegnerverhalten und Gegnerplatzierung in einem 3D-Hack-and-Slash-Prototyp konzipiert sein, damit ein zusammenhängender Kampfabschnitt mit nachvollziehbarem Spielfluss, angemessener Schwierigkeit und befriedigendem Spielgefühl entsteht?}
\end{quote}

\vspace{\baselineskip}

Diese Frage setzt sich aus zwei inhaltlichen Komponenten zusammen: Zum einen geht es um das Gegnerverhalten und das Kampfgefühl, zum anderen um die Gegnerplatzierung und das Leveldesign. Beide Komponenten werden anhand von drei Qualitäten beurteilt: Spielfluss, Schwierigkeit und Spielgefühl. Wie wir diese Begriffe genau verstehen, wird im 2. kapitel Methodik dieses Begleitberichts ersichtlich.

\begin{itemize}[leftmargin=1.5cm]
    \item \textcolor{red}{\textbf{Vertrag:} Die Erkenntnisfrage weicht vom Vertrag ab (dort: \enquote{Wie können \ldots so umgesetzt werden, dass ein Kampfabschnitt von Testspieler:innen als flüssig und herausfordernd erlebt wird?}). Mit dem Betreuer klären oder angleichen.}
    \item \textcolor{red}{\textbf{Begriffe:} \enquote{im Verlauf dieses Begleitberichts} ist zu vage, der Leser soll wissen, wo. Zurzeit definiert Kapitel~\ref{sec:Grundlagen} nur das Spielgefühl (Game Feel nach Swink, Abschnitt~2.2); Spielfluss und Schwierigkeit werden nirgends definiert. Entweder hier je ein Satz pro Begriff oder in Abschnitt~2.2 ergänzen und konkret darauf verweisen.}
    \item \textcolor{red}{\textbf{Kampfgefühl vs.\ Spielgefühl:} Wenn das zwei verschiedene Dinge sind, kurz den Unterschied erklären. Wenn nicht, ein Wort verwenden (dann auch in der Einleitung von Kapitel~\ref{sec:Umsetzung} anpassen).}
\end{itemize}

\subsection{Zielsetzung und Aufbau der Arbeit} \label{sec:Aufbau}

\textcolor{red}{Am Ende den Vertrag, das Spiel und den Text unten Abgleichen}

Ziel dieser Arbeit ist die Entwicklung eines spielbaren 3D-Hack-and-Slash-Prototyps, mit dem wir unsere Erkenntnisfrage zu beantworten versuchen.

Auf dem Weg zur Beantwortung unserer Erkenntnisfrage haben wir uns auch Teilziele gesetzt. Nähmlich die Analyse von Zwei Referenzspielen, Der Spielbare Prototyp mit einem zusammenhängendem Kampfabschnitt und die Evaluation durch Playtests.

Dafür sind wir wie folgt vorgegangen: Zuerst haben wir ausgehend von der Erkenntnisfrage zu den Themen Gegnerverhalten (Kapitel~\ref{sec:Referenzspielanalyse}) und Gegnerplatzierung (Kapitel~\ref{sec:Leveldesign}) recherchiert. Um das Gegnerverhalten zu untersuchen, haben wir Genrevertreter selbst gespielt und unsere Beobachtungen mithilfe eines eigenen Analyseframeworks (Abschnitt~\ref{sec:Analysefragebogen}) festgehalten. Für die Gegnerplatzierung haben wir Fachliteratur zum Leveldesign sowie Vorträge von Fachleuten ausgewertet. Daraus haben wir acht Designvorschläge abgeleitet, je vier zum Gegnerverhalten und zur Gegnerplatzierung, die nach unserer Einschätzung besonders viel zur Beantwortung der Erkenntnisfrage beitragen. Diese Vorschläge haben wir anschliessend konkret in unserem Prototyp umgesetzt (Kapitel~\ref{sec:Umsetzung}). In Playtests mit fünf Testpersonen haben wir dann versucht, die Designvorschläge zu evaluieren (Kapitel~\ref{sec:Evaluation}). Zum Schluss haben wir aus der Evaluation, den Playtests und unseren Recherchen das weitere Vorgehen abgeleitet und Faktoren zusammengetragen, die nach unserer Einschätzung das Gegnerverhalten und die Gegnerplatzierung ebenfalls beeinflussen und die wir als Nächstes untersuchen würden (Kapitel~\ref{sec:Schlussfolgerungen}).

Dies ist der Begleitbericht zu unserem Prototypen fussnote ergänzen bitte. Dazu gehört noch ein dokuent, das die Technische entwicklung in Godot begleitet.

\begin{itemize}[leftmargin=1.5cm]
    \item \textcolor{red}{\textbf{Ziele:} Bisher nennt ihr nur ein Ziel (den Prototyp). Analyse und Playtest stecken im Vorgehen. Die Wegleitung verlangt überprüfbare Ziele, und in Kapitel~\ref{sec:Schlussfolgerungen} prüft ihr, ob sie erreicht wurden. Vorschlag: nach dem ersten Satz die drei Teilziele als kurze Liste (Analyse zweier Referenzspiele, Prototyp mit zusammenhängendem Kampfabschnitt, Evaluation in Playtests).}
    \item \textcolor{red}{\textbf{Produkt} in zwei Sätzen: Godot, Third-Person, ein Kampfabschnitt aus zwölf Räumen, fünf Gegnertypen. Für Details auf Abschnitt~\ref{sec:Prototyp}, die technische Dokumentation und das README verweisen; Download-Link als Fussnote.}
    \item \textcolor{red}{\textbf{Abgrenzung} (1--2 Sätze): Prototyp statt fertiges Spiel; qualitative Playtests statt Statistik; technische Umsetzung in der technischen Dokumentation (Titel nennen).}
    \item \textcolor{red}{\textbf{Aufbau unvollständig:} Kapitel~\ref{sec:Grundlagen} (Grundlagen und Methodik), Kapitel~\ref{sec:Reflexion} (Reflexion) und der Anhang mit den Playtest-Protokollen (Abschnitt~\ref{sec:PlaytestProtokolle}) fehlen noch. Je ein Satz genügt.}
    \item \textcolor{red}{Kapitel~\ref{sec:Referenzspielanalyse} behandelt auch die Gegnerplatzierung (Fragen A3, D3). Ggf. \enquote{vor allem} ergänzen.}
    \item \textcolor{red}{Formatierung: Ziele und Vorgehen sind schon zwei Absätze, das passt. Den langen zweiten Absatz könnt ihr nach \enquote{\ldots im Prototyp umgesetzt} noch einmal teilen.}
\end{itemize}

\newpage

\end{document}
